% Generado por technical/generar_apendices_catalogo.py. No editar a mano.
\clearpage
\section{Catalogues de multisections, types de route et classes sectorielles}
\label{sec:vi-catalogos-estructurales}

Cette annexe organise trois niveaux distincts de l’inventaire interne. Les
13 multisections réunissent des cellules de l’atlas ; les 56 collections classent des routes
internes ; les 23 classes documentent l’articulation avec des secteurs et des
représentations physiques. Ces cardinaux ne dénombrent pas le même objet et
ne s’additionnent pas pour obtenir un nombre de particules. Les tableaux accompagnent les
définitions et les démonstrations de l’article : leur fonction est de permettre de localiser
chaque registre, ses incidences et ses lecteurs, non de remplacer une preuve par un recensement.

L’ordre documentaire des lignes, leurs identifiants et toutes
les répétitions figurant dans les sources sont conservés. Les identifiants littéraux
permettent de comparer les tableaux aux fichiers JSON. Un tiret indique une cellule
vide de la source ; il ne représente pas un zéro. Les fichiers conservent en outre tous
les champs et chiffres, les états originaux et la provenance vérifiable.
\CatalogRouteReference
% La remisión depende del punto de entrada: artículo o catálogo.

\subsection{Les 13 multisections de l’atlas}

Pour une multisection \(M\), soit \(\mathcal C(M)\) son ensemble de cellules et
soit \(\gamma_{B_1}(c)\) la route primaire que l’atlas conserve dans la cellule
\(c\). Dans cette annexe structurelle, on définit la fibre
\[
 B_1(M):=
 \operatorname{span}\bigl\{e_{\gamma_{B_1}(c)}:c\in\mathcal C(M)\bigr\}
 \subseteq \ell^2(\mathcal R_{\mathrm{adm}}).
\]
Les routes primaires de cellules distinctes déterminent des vecteurs de base distincts ;
par conséquent,
\(\operatorname{rank}B_1(M)=\dim B_1(M)=\#\mathcal C(M)\).
Cette définition fixe l’objet auquel se réfère le rang du tableau ; la
colonne des collections dénombre les
configurations de route présentes, non des espèces physiques. Les spectres entiers
des lecteurs \(K_D\) et \(K_\Omega\) sont conservés séparément.
Dans la notation littérale \texttt{x} indique une multiplicité. Les 13 lignes ont
le statut documentaire d’opérateur diagonal exact sur la fibre \(B_1\).
Sa représentation comme scalaire physique nécessite le transport déclaré, sauf
la sélection de rang un prévue dans le registre lui-même. Cet énoncé
reproduit la portée enregistrée et ne transforme pas tous les spectres en masses uniques.

\begingroup
\setlength{\tabcolsep}{4pt}
\setlength{\extrarowheight}{4pt}
\renewcommand{\arraystretch}{1.13}
\begin{longtable}{@{}>{\raggedright\arraybackslash}p{\dimexpr 0.19871795\linewidth-2\tabcolsep\relax}>{\raggedright\arraybackslash}p{\dimexpr 0.19871795\linewidth-2\tabcolsep\relax}>{\raggedright\arraybackslash}p{\dimexpr 0.30769231\linewidth-2\tabcolsep\relax}>{\raggedright\arraybackslash}p{\dimexpr 0.29487179\linewidth-2\tabcolsep\relax}@{}}
\caption{Les 13 multisections, sans sélection de lignes.}\label{tab:vi-multisecciones}\\
\toprule
Multisection & Rang et incidence & Classification et profondeur & Lecteurs entiers \\
\midrule
\endfirsthead
\multicolumn{4}{l}{\textit{Suite}}\\
\toprule
Multisection & Rang et incidence & Classification et profondeur & Lecteurs entiers \\
\midrule
\endhead
\midrule
\multicolumn{4}{r}{Suite à la page suivante}\\
\endfoot
\bottomrule
\endlastfoot
% VI-CATALOG-ROW multisecciones 1
charged\_\allowbreak{}lepton\_\allowbreak{}G0 & \textit{Rang}: 1\par \textit{Cellules (identifiants)}: 41\par \textit{Collections}: 1 & \textit{Collection}: charged\_\allowbreak{}lepton\par \textit{Profondeur}: G0\_\allowbreak{}center\par \textit{Couleur}: none & \textit{K D}: (80,\allowbreak{}54,\allowbreak{}6)\allowbreak{}x1\par \textit{K Ω}: (80,\allowbreak{}54,\allowbreak{}6)\allowbreak{}x1 \\
% VI-CATALOG-ROW multisecciones 2
charged\_\allowbreak{}lepton\_\allowbreak{}G2 & \textit{Rang}: 8\par \textit{Cellules (identifiants)}: 21,\allowbreak{}23,\allowbreak{}25,\allowbreak{}39,\allowbreak{}43,\allowbreak{}57,\allowbreak{}59,\allowbreak{}61\par \textit{Collections}: 8 & \textit{Collection}: charged\_\allowbreak{}lepton\par \textit{Profondeur}: G2\par \textit{Couleur}: none & \textit{K D}: (0,\allowbreak{}24,\allowbreak{}0)\allowbreak{}x1; (40,\allowbreak{}78,\allowbreak{}27)\allowbreak{}x2; (80,\allowbreak{}54,\allowbreak{}7)\allowbreak{}x1; (80,\allowbreak{}66,\allowbreak{}2)\allowbreak{}x1; (80,\allowbreak{}66,\allowbreak{}3)\allowbreak{}x1; (100,\allowbreak{}66,\allowbreak{}0)\allowbreak{}x1; (100,\allowbreak{}66,\allowbreak{}2)\allowbreak{}x1\par \textit{K Ω}: (80,\allowbreak{}54,\allowbreak{}6)\allowbreak{}x8 \\
% VI-CATALOG-ROW multisecciones 3
charged\_\allowbreak{}lepton\_\allowbreak{}G4 & \textit{Rang}: 16\par \textit{Cellules (identifiants)}: 1,\allowbreak{}3,\allowbreak{}5,\allowbreak{}7,\allowbreak{}9,\allowbreak{}19,\allowbreak{}27,\allowbreak{}37,\allowbreak{}45,\allowbreak{}55,\allowbreak{}63,\allowbreak{}73,\allowbreak{}75,\allowbreak{}77,\allowbreak{}79,\allowbreak{}81\par \textit{Collections}: 16 & \textit{Collection}: charged\_\allowbreak{}lepton\par \textit{Profondeur}: G4\par \textit{Couleur}: none & \textit{K D}: (0,\allowbreak{}24,\allowbreak{}0)\allowbreak{}x1; (40,\allowbreak{}24,\allowbreak{}2)\allowbreak{}x2; (40,\allowbreak{}54,\allowbreak{}6)\allowbreak{}x2; (40,\allowbreak{}84,\allowbreak{}30)\allowbreak{}x2; (60,\allowbreak{}78,\allowbreak{}25)\allowbreak{}x3; (80,\allowbreak{}66,\allowbreak{}2)\allowbreak{}x2; (80,\allowbreak{}66,\allowbreak{}3)\allowbreak{}x3; (80,\allowbreak{}66,\allowbreak{}4)\allowbreak{}x1\par \textit{K Ω}: (24,\allowbreak{}54,\allowbreak{}2)\allowbreak{}x1; (56,\allowbreak{}54,\allowbreak{}5)\allowbreak{}x4; (80,\allowbreak{}54,\allowbreak{}6)\allowbreak{}x11 \\
% VI-CATALOG-ROW multisecciones 4
neutrino\_\allowbreak{}G1 & \textit{Rang}: 2\par \textit{Cellules (identifiants)}: 40,\allowbreak{}42\par \textit{Collections}: 2 & \textit{Collection}: neutrino\par \textit{Profondeur}: G1\par \textit{Couleur}: none & \textit{K D}: (40,\allowbreak{}24,\allowbreak{}2)\allowbreak{}x1; (40,\allowbreak{}78,\allowbreak{}27)\allowbreak{}x1\par \textit{K Ω}: (40,\allowbreak{}54,\allowbreak{}4)\allowbreak{}x1; (80,\allowbreak{}54,\allowbreak{}6)\allowbreak{}x1 \\
% VI-CATALOG-ROW multisecciones 5
neutrino\_\allowbreak{}G2 & \textit{Rang}: 4\par \textit{Cellules (identifiants)}: 22,\allowbreak{}24,\allowbreak{}58,\allowbreak{}60\par \textit{Collections}: 4 & \textit{Collection}: neutrino\par \textit{Profondeur}: G2\par \textit{Couleur}: none & \textit{K D}: (0,\allowbreak{}24,\allowbreak{}0)\allowbreak{}x1; (40,\allowbreak{}84,\allowbreak{}30)\allowbreak{}x1; (60,\allowbreak{}78,\allowbreak{}25)\allowbreak{}x1; (80,\allowbreak{}66,\allowbreak{}3)\allowbreak{}x1\par \textit{K Ω}: (4,\allowbreak{}54,\allowbreak{}2)\allowbreak{}x1; (80,\allowbreak{}54,\allowbreak{}6)\allowbreak{}x3 \\
% VI-CATALOG-ROW multisecciones 6
neutrino\_\allowbreak{}G3 & \textit{Rang}: 6\par \textit{Cellules (identifiants)}: 20,\allowbreak{}26,\allowbreak{}38,\allowbreak{}44,\allowbreak{}56,\allowbreak{}62\par \textit{Collections}: 6 & \textit{Collection}: neutrino\par \textit{Profondeur}: G3\par \textit{Couleur}: none & \textit{K D}: (40,\allowbreak{}24,\allowbreak{}2)\allowbreak{}x2; (40,\allowbreak{}78,\allowbreak{}27)\allowbreak{}x1; (60,\allowbreak{}84,\allowbreak{}28)\allowbreak{}x1; (80,\allowbreak{}54,\allowbreak{}6)\allowbreak{}x1; (80,\allowbreak{}66,\allowbreak{}3)\allowbreak{}x1\par \textit{K Ω}: (36,\allowbreak{}54,\allowbreak{}4)\allowbreak{}x1; (56,\allowbreak{}54,\allowbreak{}5)\allowbreak{}x1; (80,\allowbreak{}54,\allowbreak{}6)\allowbreak{}x4 \\
% VI-CATALOG-ROW multisecciones 7
neutrino\_\allowbreak{}G4 & \textit{Rang}: 8\par \textit{Cellules (identifiants)}: 2,\allowbreak{}4,\allowbreak{}6,\allowbreak{}8,\allowbreak{}74,\allowbreak{}76,\allowbreak{}78,\allowbreak{}80\par \textit{Collections}: 8 & \textit{Collection}: neutrino\par \textit{Profondeur}: G4\par \textit{Couleur}: none & \textit{K D}: (0,\allowbreak{}24,\allowbreak{}0)\allowbreak{}x1; (40,\allowbreak{}24,\allowbreak{}2)\allowbreak{}x1; (40,\allowbreak{}78,\allowbreak{}27)\allowbreak{}x2; (40,\allowbreak{}84,\allowbreak{}30)\allowbreak{}x1; (80,\allowbreak{}54,\allowbreak{}7)\allowbreak{}x1; (80,\allowbreak{}66,\allowbreak{}2)\allowbreak{}x1; (100,\allowbreak{}66,\allowbreak{}0)\allowbreak{}x1\par \textit{K Ω}: (36,\allowbreak{}54,\allowbreak{}4)\allowbreak{}x1; (56,\allowbreak{}54,\allowbreak{}5)\allowbreak{}x1; (80,\allowbreak{}54,\allowbreak{}6)\allowbreak{}x6 \\
% VI-CATALOG-ROW multisecciones 8
quark\_\allowbreak{}down\_\allowbreak{}G1 & \textit{Rang}: 2\par \textit{Cellules (identifiants)}: 32,\allowbreak{}50\par \textit{Collections}: 2 & \textit{Collection}: quark\_\allowbreak{}down\par \textit{Profondeur}: G1\par \textit{Couleur}: B | G & \textit{K D}: (40,\allowbreak{}24,\allowbreak{}2)\allowbreak{}x1; (40,\allowbreak{}78,\allowbreak{}27)\allowbreak{}x1\par \textit{K Ω}: (40,\allowbreak{}54,\allowbreak{}4)\allowbreak{}x1; (80,\allowbreak{}54,\allowbreak{}6)\allowbreak{}x1 \\
% VI-CATALOG-ROW multisecciones 9
quark\_\allowbreak{}down\_\allowbreak{}G2 & \textit{Rang}: 4\par \textit{Cellules (identifiants)}: 30,\allowbreak{}34,\allowbreak{}48,\allowbreak{}52\par \textit{Collections}: 4 & \textit{Collection}: quark\_\allowbreak{}down\par \textit{Profondeur}: G2\par \textit{Couleur}: B | G | R & \textit{K D}: (0,\allowbreak{}24,\allowbreak{}0)\allowbreak{}x1; (40,\allowbreak{}84,\allowbreak{}30)\allowbreak{}x1; (60,\allowbreak{}78,\allowbreak{}25)\allowbreak{}x1; (80,\allowbreak{}66,\allowbreak{}2)\allowbreak{}x1\par \textit{K Ω}: (4,\allowbreak{}54,\allowbreak{}2)\allowbreak{}x1; (80,\allowbreak{}54,\allowbreak{}6)\allowbreak{}x3 \\
% VI-CATALOG-ROW multisecciones 10
quark\_\allowbreak{}down\_\allowbreak{}G3 & \textit{Rang}: 6\par \textit{Cellules (identifiants)}: 12,\allowbreak{}14,\allowbreak{}16,\allowbreak{}66,\allowbreak{}68,\allowbreak{}70\par \textit{Collections}: 6 & \textit{Collection}: quark\_\allowbreak{}down\par \textit{Profondeur}: G3\par \textit{Couleur}: B | G | R & \textit{K D}: (40,\allowbreak{}24,\allowbreak{}2)\allowbreak{}x2; (40,\allowbreak{}78,\allowbreak{}27)\allowbreak{}x1; (60,\allowbreak{}84,\allowbreak{}28)\allowbreak{}x1; (80,\allowbreak{}54,\allowbreak{}6)\allowbreak{}x1; (80,\allowbreak{}66,\allowbreak{}4)\allowbreak{}x1\par \textit{K Ω}: (36,\allowbreak{}54,\allowbreak{}4)\allowbreak{}x1; (80,\allowbreak{}54,\allowbreak{}6)\allowbreak{}x5 \\
% VI-CATALOG-ROW multisecciones 11
quark\_\allowbreak{}down\_\allowbreak{}G4 & \textit{Rang}: 8\par \textit{Cellules (identifiants)}: 10,\allowbreak{}18,\allowbreak{}28,\allowbreak{}36,\allowbreak{}46,\allowbreak{}54,\allowbreak{}64,\allowbreak{}72\par \textit{Collections}: 8 & \textit{Collection}: quark\_\allowbreak{}down\par \textit{Profondeur}: G4\par \textit{Couleur}: B | G | R & \textit{K D}: (0,\allowbreak{}24,\allowbreak{}0)\allowbreak{}x1; (40,\allowbreak{}24,\allowbreak{}2)\allowbreak{}x1; (40,\allowbreak{}78,\allowbreak{}27)\allowbreak{}x2; (40,\allowbreak{}84,\allowbreak{}30)\allowbreak{}x1; (80,\allowbreak{}54,\allowbreak{}7)\allowbreak{}x1; (80,\allowbreak{}66,\allowbreak{}3)\allowbreak{}x1; (100,\allowbreak{}66,\allowbreak{}2)\allowbreak{}x1\par \textit{K Ω}: (36,\allowbreak{}54,\allowbreak{}4)\allowbreak{}x1; (80,\allowbreak{}54,\allowbreak{}6)\allowbreak{}x7 \\
% VI-CATALOG-ROW multisecciones 12
quark\_\allowbreak{}up\_\allowbreak{}G1 & \textit{Rang}: 4\par \textit{Cellules (identifiants)}: 31,\allowbreak{}33,\allowbreak{}49,\allowbreak{}51\par \textit{Collections}: 3 & \textit{Collection}: quark\_\allowbreak{}up\par \textit{Profondeur}: G1\par \textit{Couleur}: B | G | R & \textit{K D}: (40,\allowbreak{}54,\allowbreak{}6)\allowbreak{}x1; (60,\allowbreak{}78,\allowbreak{}25)\allowbreak{}x1; (80,\allowbreak{}66,\allowbreak{}2)\allowbreak{}x1; (80,\allowbreak{}66,\allowbreak{}3)\allowbreak{}x1\par \textit{K Ω}: (80,\allowbreak{}54,\allowbreak{}6)\allowbreak{}x4 \\
% VI-CATALOG-ROW multisecciones 13
quark\_\allowbreak{}up\_\allowbreak{}G3 & \textit{Rang}: 12\par \textit{Cellules (identifiants)}: 11,\allowbreak{}13,\allowbreak{}15,\allowbreak{}17,\allowbreak{}29,\allowbreak{}35,\allowbreak{}47,\allowbreak{}53,\allowbreak{}65,\allowbreak{}67,\allowbreak{}69,\allowbreak{}71\par \textit{Collections}: 12 & \textit{Collection}: quark\_\allowbreak{}up\par \textit{Profondeur}: G3\par \textit{Couleur}: B | G | R & \textit{K D}: (40,\allowbreak{}24,\allowbreak{}2)\allowbreak{}x2; (40,\allowbreak{}78,\allowbreak{}27)\allowbreak{}x2; (60,\allowbreak{}84,\allowbreak{}28)\allowbreak{}x1; (80,\allowbreak{}54,\allowbreak{}6)\allowbreak{}x3; (80,\allowbreak{}66,\allowbreak{}3)\allowbreak{}x1; (80,\allowbreak{}66,\allowbreak{}4)\allowbreak{}x1; (100,\allowbreak{}66,\allowbreak{}0)\allowbreak{}x1; (100,\allowbreak{}66,\allowbreak{}2)\allowbreak{}x1\par \textit{K Ω}: (56,\allowbreak{}54,\allowbreak{}5)\allowbreak{}x3; (80,\allowbreak{}54,\allowbreak{}6)\allowbreak{}x9 \\
\end{longtable}
\endgroup
\subsection{Classification des 56 types de route}

La signature qui identifie une collection est conservée littéralement, avec ses signes
et ses séparateurs. Les cellules et les multisections rendent visible qu’une collection de
route n’équivaut pas par définition à une espèce physique. Les profils conjoints
dénombrent les combinaisons enregistrées dans la collection. Les spectres
décimaux du caractère et les spectres complets de signature restent intacts
dans \texttt{catalogo\_familias.json} ; ici sont exprimés l’incidence et les
lecteurs entiers qui permettent de s’orienter dans cette information.

\begingroup
\setlength{\tabcolsep}{4pt}
\setlength{\extrarowheight}{4pt}
\renewcommand{\arraystretch}{1.13}
\begin{longtable}{@{}>{\raggedright\arraybackslash}p{\dimexpr 0.21794872\linewidth-2\tabcolsep\relax}>{\raggedright\arraybackslash}p{\dimexpr 0.34615385\linewidth-2\tabcolsep\relax}>{\raggedright\arraybackslash}p{\dimexpr 0.21794872\linewidth-2\tabcolsep\relax}>{\raggedright\arraybackslash}p{\dimexpr 0.21794872\linewidth-2\tabcolsep\relax}@{}}
\caption{Les 56 collections internes de route.}\label{tab:vi-familias}\\
\toprule
Collection et signature & Incidence & \(K_D\) & \(K_\Omega\) \\
\midrule
\endfirsthead
\multicolumn{4}{l}{\textit{Suite}}\\
\toprule
Collection et signature & Incidence & \(K_D\) & \(K_\Omega\) \\
\midrule
\endhead
\midrule
\multicolumn{4}{r}{Suite à la page suivante}\\
\endfoot
\bottomrule
\endlastfoot
% VI-CATALOG-ROW familias 1
1\par 24|\allowbreak{}0|\allowbreak{}12|\allowbreak{}0|\allowbreak{}-12|\allowbreak{}+++-\kern0pt{}-\kern0pt{}-+++-\kern0pt{}-\kern0pt{}- & \textit{Nombre de cellules}: 5\par \textit{Cellules (identifiants)}: 21,\allowbreak{}31,\allowbreak{}41,\allowbreak{}51,\allowbreak{}81\par \textit{Multisections}: charged\_\allowbreak{}lepton\_\allowbreak{}G0 | charged\_\allowbreak{}lepton\_\allowbreak{}G2 | charged\_\allowbreak{}lepton\_\allowbreak{}G4 | quark\_\allowbreak{}up\_\allowbreak{}G1\par \textit{Collections}: charged\_\allowbreak{}lepton | quark\_\allowbreak{}up\par \textit{Couleur}: R | none\par \textit{Profils conjoints}: 5 & (0,\allowbreak{}24,\allowbreak{}0)\allowbreak{}x2; (40,\allowbreak{}54,\allowbreak{}6)\allowbreak{}x1; (60,\allowbreak{}78,\allowbreak{}25)\allowbreak{}x1; (80,\allowbreak{}54,\allowbreak{}6)\allowbreak{}x1 & (24,\allowbreak{}54,\allowbreak{}2)\allowbreak{}x1; (80,\allowbreak{}54,\allowbreak{}6)\allowbreak{}x4 \\
% VI-CATALOG-ROW familias 2
2\par 24|\allowbreak{}0|\allowbreak{}4|\allowbreak{}0|\allowbreak{}-12|\allowbreak{}+++-\kern0pt{}-\kern0pt{}-+++-\kern0pt{}-\kern0pt{}- & \textit{Nombre de cellules}: 2\par \textit{Cellules (identifiants)}: 11,\allowbreak{}61\par \textit{Multisections}: charged\_\allowbreak{}lepton\_\allowbreak{}G2 | quark\_\allowbreak{}up\_\allowbreak{}G3\par \textit{Collections}: charged\_\allowbreak{}lepton | quark\_\allowbreak{}up\par \textit{Couleur}: R | none\par \textit{Profils conjoints}: 2 & (60,\allowbreak{}84,\allowbreak{}28)\allowbreak{}x1; (80,\allowbreak{}54,\allowbreak{}7)\allowbreak{}x1 & (80,\allowbreak{}54,\allowbreak{}6)\allowbreak{}x2 \\
% VI-CATALOG-ROW familias 3
3\par 24|\allowbreak{}0|\allowbreak{}8|\allowbreak{}0|\allowbreak{}-12|\allowbreak{}+++-\kern0pt{}-\kern0pt{}-+++-\kern0pt{}-\kern0pt{}- & \textit{Nombre de cellules}: 2\par \textit{Cellules (identifiants)}: 1,\allowbreak{}71\par \textit{Multisections}: charged\_\allowbreak{}lepton\_\allowbreak{}G4 | quark\_\allowbreak{}up\_\allowbreak{}G3\par \textit{Collections}: charged\_\allowbreak{}lepton | quark\_\allowbreak{}up\par \textit{Couleur}: R | none\par \textit{Profils conjoints}: 2 & (60,\allowbreak{}78,\allowbreak{}25)\allowbreak{}x1; (80,\allowbreak{}54,\allowbreak{}6)\allowbreak{}x1 & (80,\allowbreak{}54,\allowbreak{}6)\allowbreak{}x2 \\
% VI-CATALOG-ROW familias 4
4\par 28|\allowbreak{}-6|\allowbreak{}-2|\allowbreak{}-2|\allowbreak{}-8|\allowbreak{}0++-\kern0pt{}-\kern0pt{}-0++-\kern0pt{}-\kern0pt{}- & \textit{Nombre de cellules}: 1\par \textit{Cellules (identifiants)}: 53\par \textit{Multisections}: quark\_\allowbreak{}up\_\allowbreak{}G3\par \textit{Collections}: quark\_\allowbreak{}up\par \textit{Couleur}: G\par \textit{Profils conjoints}: 1 & (80,\allowbreak{}66,\allowbreak{}3)\allowbreak{}x1 & (56,\allowbreak{}54,\allowbreak{}5)\allowbreak{}x1 \\
% VI-CATALOG-ROW familias 5
5\par 28|\allowbreak{}-6|\allowbreak{}-4|\allowbreak{}-2|\allowbreak{}-8|\allowbreak{}0++-\kern0pt{}-\kern0pt{}-0++-\kern0pt{}-\kern0pt{}- & \textit{Nombre de cellules}: 1\par \textit{Cellules (identifiants)}: 20\par \textit{Multisections}: neutrino\_\allowbreak{}G3\par \textit{Collections}: neutrino\par \textit{Couleur}: none\par \textit{Profils conjoints}: 1 & (80,\allowbreak{}66,\allowbreak{}3)\allowbreak{}x1 & (56,\allowbreak{}54,\allowbreak{}5)\allowbreak{}x1 \\
% VI-CATALOG-ROW familias 6
6\par 28|\allowbreak{}-6|\allowbreak{}0|\allowbreak{}0|\allowbreak{}-12|\allowbreak{}+++-\kern0pt{}-\kern0pt{}-+++-\kern0pt{}-\kern0pt{}- & \textit{Nombre de cellules}: 2\par \textit{Cellules (identifiants)}: 30,\allowbreak{}70\par \textit{Multisections}: quark\_\allowbreak{}down\_\allowbreak{}G2 | quark\_\allowbreak{}down\_\allowbreak{}G3\par \textit{Collections}: quark\_\allowbreak{}down\par \textit{Couleur}: G\par \textit{Profils conjoints}: 2 & (40,\allowbreak{}24,\allowbreak{}2)\allowbreak{}x1; (80,\allowbreak{}66,\allowbreak{}2)\allowbreak{}x1 & (80,\allowbreak{}54,\allowbreak{}6)\allowbreak{}x2 \\
% VI-CATALOG-ROW familias 7
7\par 28|\allowbreak{}-6|\allowbreak{}10|\allowbreak{}0|\allowbreak{}-12|\allowbreak{}+++-\kern0pt{}-\kern0pt{}-+++-\kern0pt{}-\kern0pt{}- & \textit{Nombre de cellules}: 2\par \textit{Cellules (identifiants)}: 33,\allowbreak{}64\par \textit{Multisections}: quark\_\allowbreak{}down\_\allowbreak{}G4 | quark\_\allowbreak{}up\_\allowbreak{}G1\par \textit{Collections}: quark\_\allowbreak{}down | quark\_\allowbreak{}up\par \textit{Couleur}: G\par \textit{Profils conjoints}: 2 & (40,\allowbreak{}24,\allowbreak{}2)\allowbreak{}x1; (80,\allowbreak{}66,\allowbreak{}2)\allowbreak{}x1 & (36,\allowbreak{}54,\allowbreak{}4)\allowbreak{}x1; (80,\allowbreak{}54,\allowbreak{}6)\allowbreak{}x1 \\
% VI-CATALOG-ROW familias 8
8\par 28|\allowbreak{}-6|\allowbreak{}12|\allowbreak{}0|\allowbreak{}-12|\allowbreak{}+++-\kern0pt{}-\kern0pt{}-+++-\kern0pt{}-\kern0pt{}- & \textit{Nombre de cellules}: 1\par \textit{Cellules (identifiants)}: 10\par \textit{Multisections}: quark\_\allowbreak{}down\_\allowbreak{}G4\par \textit{Collections}: quark\_\allowbreak{}down\par \textit{Couleur}: G\par \textit{Profils conjoints}: 1 & (100,\allowbreak{}66,\allowbreak{}2)\allowbreak{}x1 & (80,\allowbreak{}54,\allowbreak{}6)\allowbreak{}x1 \\
% VI-CATALOG-ROW familias 9
9\par 28|\allowbreak{}-6|\allowbreak{}2|\allowbreak{}0|\allowbreak{}-12|\allowbreak{}+++-\kern0pt{}-\kern0pt{}-+++-\kern0pt{}-\kern0pt{}- & \textit{Nombre de cellules}: 1\par \textit{Cellules (identifiants)}: 3\par \textit{Multisections}: charged\_\allowbreak{}lepton\_\allowbreak{}G4\par \textit{Collections}: charged\_\allowbreak{}lepton\par \textit{Couleur}: none\par \textit{Profils conjoints}: 1 & (40,\allowbreak{}84,\allowbreak{}30)\allowbreak{}x1 & (80,\allowbreak{}54,\allowbreak{}6)\allowbreak{}x1 \\
% VI-CATALOG-ROW familias 10
10\par 28|\allowbreak{}-6|\allowbreak{}4|\allowbreak{}-2|\allowbreak{}-8|\allowbreak{}++0-\kern0pt{}-\kern0pt{}-++0-\kern0pt{}-\kern0pt{}- & \textit{Nombre de cellules}: 1\par \textit{Cellules (identifiants)}: 50\par \textit{Multisections}: quark\_\allowbreak{}down\_\allowbreak{}G1\par \textit{Collections}: quark\_\allowbreak{}down\par \textit{Couleur}: G\par \textit{Profils conjoints}: 1 & (40,\allowbreak{}78,\allowbreak{}27)\allowbreak{}x1 & (80,\allowbreak{}54,\allowbreak{}6)\allowbreak{}x1 \\
% VI-CATALOG-ROW familias 11
11\par 28|\allowbreak{}-6|\allowbreak{}4|\allowbreak{}0|\allowbreak{}-12|\allowbreak{}+++-\kern0pt{}-\kern0pt{}-+++-\kern0pt{}-\kern0pt{}- & \textit{Nombre de cellules}: 2\par \textit{Cellules (identifiants)}: 9,\allowbreak{}40\par \textit{Multisections}: charged\_\allowbreak{}lepton\_\allowbreak{}G4 | neutrino\_\allowbreak{}G1\par \textit{Collections}: charged\_\allowbreak{}lepton | neutrino\par \textit{Couleur}: none\par \textit{Profils conjoints}: 2 & (40,\allowbreak{}24,\allowbreak{}2)\allowbreak{}x1; (80,\allowbreak{}66,\allowbreak{}2)\allowbreak{}x1 & (40,\allowbreak{}54,\allowbreak{}4)\allowbreak{}x1; (80,\allowbreak{}54,\allowbreak{}6)\allowbreak{}x1 \\
% VI-CATALOG-ROW familias 12
12\par 28|\allowbreak{}-6|\allowbreak{}4|\allowbreak{}0|\allowbreak{}-8|\allowbreak{}+0+-\kern0pt{}-\kern0pt{}-+0+-\kern0pt{}-\kern0pt{}- & \textit{Nombre de cellules}: 1\par \textit{Cellules (identifiants)}: 80\par \textit{Multisections}: neutrino\_\allowbreak{}G4\par \textit{Collections}: neutrino\par \textit{Couleur}: none\par \textit{Profils conjoints}: 1 & (40,\allowbreak{}78,\allowbreak{}27)\allowbreak{}x1 & (80,\allowbreak{}54,\allowbreak{}6)\allowbreak{}x1 \\
% VI-CATALOG-ROW familias 13
13\par 28|\allowbreak{}-6|\allowbreak{}6|\allowbreak{}-2|\allowbreak{}-8|\allowbreak{}++0-\kern0pt{}-\kern0pt{}-++0-\kern0pt{}-\kern0pt{}- & \textit{Nombre de cellules}: 1\par \textit{Cellules (identifiants)}: 74\par \textit{Multisections}: neutrino\_\allowbreak{}G4\par \textit{Collections}: neutrino\par \textit{Couleur}: none\par \textit{Profils conjoints}: 1 & (40,\allowbreak{}78,\allowbreak{}27)\allowbreak{}x1 & (80,\allowbreak{}54,\allowbreak{}6)\allowbreak{}x1 \\
% VI-CATALOG-ROW familias 14
14\par 28|\allowbreak{}-6|\allowbreak{}6|\allowbreak{}0|\allowbreak{}-12|\allowbreak{}+++-\kern0pt{}-\kern0pt{}-+++-\kern0pt{}-\kern0pt{}- & \textit{Nombre de cellules}: 3\par \textit{Cellules (identifiants)}: 13,\allowbreak{}43,\allowbreak{}63\par \textit{Multisections}: charged\_\allowbreak{}lepton\_\allowbreak{}G2 | charged\_\allowbreak{}lepton\_\allowbreak{}G4 | quark\_\allowbreak{}up\_\allowbreak{}G3\par \textit{Collections}: charged\_\allowbreak{}lepton | quark\_\allowbreak{}up\par \textit{Couleur}: G | none\par \textit{Profils conjoints}: 3 & (40,\allowbreak{}24,\allowbreak{}2)\allowbreak{}x1; (80,\allowbreak{}66,\allowbreak{}2)\allowbreak{}x1; (100,\allowbreak{}66,\allowbreak{}2)\allowbreak{}x1 & (80,\allowbreak{}54,\allowbreak{}6)\allowbreak{}x3 \\
% VI-CATALOG-ROW familias 15
15\par 28|\allowbreak{}-6|\allowbreak{}6|\allowbreak{}0|\allowbreak{}-8|\allowbreak{}+0+-\kern0pt{}-\kern0pt{}-+0+-\kern0pt{}-\kern0pt{}- & \textit{Nombre de cellules}: 1\par \textit{Cellules (identifiants)}: 23\par \textit{Multisections}: charged\_\allowbreak{}lepton\_\allowbreak{}G2\par \textit{Collections}: charged\_\allowbreak{}lepton\par \textit{Couleur}: none\par \textit{Profils conjoints}: 1 & (40,\allowbreak{}78,\allowbreak{}27)\allowbreak{}x1 & (80,\allowbreak{}54,\allowbreak{}6)\allowbreak{}x1 \\
% VI-CATALOG-ROW familias 16
16\par 28|\allowbreak{}-6|\allowbreak{}8|\allowbreak{}0|\allowbreak{}-12|\allowbreak{}+++-\kern0pt{}-\kern0pt{}-+++-\kern0pt{}-\kern0pt{}- & \textit{Nombre de cellules}: 1\par \textit{Cellules (identifiants)}: 60\par \textit{Multisections}: neutrino\_\allowbreak{}G2\par \textit{Collections}: neutrino\par \textit{Couleur}: none\par \textit{Profils conjoints}: 1 & (40,\allowbreak{}84,\allowbreak{}30)\allowbreak{}x1 & (80,\allowbreak{}54,\allowbreak{}6)\allowbreak{}x1 \\
% VI-CATALOG-ROW familias 17
17\par 28|\allowbreak{}6|\allowbreak{}-10|\allowbreak{}0|\allowbreak{}-12|\allowbreak{}+++-\kern0pt{}-\kern0pt{}-+++-\kern0pt{}-\kern0pt{}- & \textit{Nombre de cellules}: 1\par \textit{Cellules (identifiants)}: 19\par \textit{Multisections}: charged\_\allowbreak{}lepton\_\allowbreak{}G4\par \textit{Collections}: charged\_\allowbreak{}lepton\par \textit{Couleur}: none\par \textit{Profils conjoints}: 1 & (40,\allowbreak{}84,\allowbreak{}30)\allowbreak{}x1 & (80,\allowbreak{}54,\allowbreak{}6)\allowbreak{}x1 \\
% VI-CATALOG-ROW familias 18
18\par 28|\allowbreak{}6|\allowbreak{}-2|\allowbreak{}0|\allowbreak{}-12|\allowbreak{}+++-\kern0pt{}-\kern0pt{}-+++-\kern0pt{}-\kern0pt{}- & \textit{Nombre de cellules}: 5\par \textit{Cellules (identifiants)}: 18,\allowbreak{}29,\allowbreak{}49,\allowbreak{}59,\allowbreak{}79\par \textit{Multisections}: charged\_\allowbreak{}lepton\_\allowbreak{}G2 | charged\_\allowbreak{}lepton\_\allowbreak{}G4 | quark\_\allowbreak{}down\_\allowbreak{}G4 | quark\_\allowbreak{}up\_\allowbreak{}G1 | quark\_\allowbreak{}up\_\allowbreak{}G3\par \textit{Collections}: charged\_\allowbreak{}lepton | quark\_\allowbreak{}down | quark\_\allowbreak{}up\par \textit{Couleur}: B | none\par \textit{Profils conjoints}: 4 & (40,\allowbreak{}24,\allowbreak{}2)\allowbreak{}x1; (40,\allowbreak{}78,\allowbreak{}27)\allowbreak{}x1; (80,\allowbreak{}66,\allowbreak{}3)\allowbreak{}x2; (100,\allowbreak{}66,\allowbreak{}0)\allowbreak{}x1 & (80,\allowbreak{}54,\allowbreak{}6)\allowbreak{}x5 \\
% VI-CATALOG-ROW familias 19
19\par 28|\allowbreak{}6|\allowbreak{}-4|\allowbreak{}0|\allowbreak{}-12|\allowbreak{}+++-\kern0pt{}-\kern0pt{}-+++-\kern0pt{}-\kern0pt{}- & \textit{Nombre de cellules}: 2\par \textit{Cellules (identifiants)}: 2,\allowbreak{}42\par \textit{Multisections}: neutrino\_\allowbreak{}G1 | neutrino\_\allowbreak{}G4\par \textit{Collections}: neutrino\par \textit{Couleur}: none\par \textit{Profils conjoints}: 2 & (40,\allowbreak{}78,\allowbreak{}27)\allowbreak{}x1; (100,\allowbreak{}66,\allowbreak{}0)\allowbreak{}x1 & (80,\allowbreak{}54,\allowbreak{}6)\allowbreak{}x2 \\
% VI-CATALOG-ROW familias 20
20\par 28|\allowbreak{}6|\allowbreak{}-6|\allowbreak{}0|\allowbreak{}-12|\allowbreak{}+++-\kern0pt{}-\kern0pt{}-+++-\kern0pt{}-\kern0pt{}- & \textit{Nombre de cellules}: 1\par \textit{Cellules (identifiants)}: 69\par \textit{Multisections}: quark\_\allowbreak{}up\_\allowbreak{}G3\par \textit{Collections}: quark\_\allowbreak{}up\par \textit{Couleur}: B\par \textit{Profils conjoints}: 1 & (80,\allowbreak{}66,\allowbreak{}4)\allowbreak{}x1 & (80,\allowbreak{}54,\allowbreak{}6)\allowbreak{}x1 \\
% VI-CATALOG-ROW familias 21
21\par 28|\allowbreak{}6|\allowbreak{}-8|\allowbreak{}0|\allowbreak{}-12|\allowbreak{}+++-\kern0pt{}-\kern0pt{}-+++-\kern0pt{}-\kern0pt{}- & \textit{Nombre de cellules}: 1\par \textit{Cellules (identifiants)}: 52\par \textit{Multisections}: quark\_\allowbreak{}down\_\allowbreak{}G2\par \textit{Collections}: quark\_\allowbreak{}down\par \textit{Couleur}: B\par \textit{Profils conjoints}: 1 & (40,\allowbreak{}84,\allowbreak{}30)\allowbreak{}x1 & (80,\allowbreak{}54,\allowbreak{}6)\allowbreak{}x1 \\
% VI-CATALOG-ROW familias 22
22\par 28|\allowbreak{}6|\allowbreak{}0|\allowbreak{}0|\allowbreak{}-12|\allowbreak{}+++-\kern0pt{}-\kern0pt{}-+++-\kern0pt{}-\kern0pt{}- & \textit{Nombre de cellules}: 4\par \textit{Cellules (identifiants)}: 22,\allowbreak{}32,\allowbreak{}72,\allowbreak{}73\par \textit{Multisections}: charged\_\allowbreak{}lepton\_\allowbreak{}G4 | neutrino\_\allowbreak{}G2 | quark\_\allowbreak{}down\_\allowbreak{}G1 | quark\_\allowbreak{}down\_\allowbreak{}G4\par \textit{Collections}: charged\_\allowbreak{}lepton | neutrino | quark\_\allowbreak{}down\par \textit{Couleur}: B | none\par \textit{Profils conjoints}: 3 & (40,\allowbreak{}24,\allowbreak{}2)\allowbreak{}x1; (40,\allowbreak{}78,\allowbreak{}27)\allowbreak{}x1; (80,\allowbreak{}66,\allowbreak{}3)\allowbreak{}x2 & (40,\allowbreak{}54,\allowbreak{}4)\allowbreak{}x1; (80,\allowbreak{}54,\allowbreak{}6)\allowbreak{}x3 \\
% VI-CATALOG-ROW familias 23
23\par 28|\allowbreak{}6|\allowbreak{}2|\allowbreak{}0|\allowbreak{}-12|\allowbreak{}+++-\kern0pt{}-\kern0pt{}-+++-\kern0pt{}-\kern0pt{}- & \textit{Nombre de cellules}: 2\par \textit{Cellules (identifiants)}: 8,\allowbreak{}39\par \textit{Multisections}: charged\_\allowbreak{}lepton\_\allowbreak{}G2 | neutrino\_\allowbreak{}G4\par \textit{Collections}: charged\_\allowbreak{}lepton | neutrino\par \textit{Couleur}: none\par \textit{Profils conjoints}: 2 & (40,\allowbreak{}24,\allowbreak{}2)\allowbreak{}x1; (40,\allowbreak{}78,\allowbreak{}27)\allowbreak{}x1 & (36,\allowbreak{}54,\allowbreak{}4)\allowbreak{}x1; (80,\allowbreak{}54,\allowbreak{}6)\allowbreak{}x1 \\
% VI-CATALOG-ROW familias 24
24\par 28|\allowbreak{}6|\allowbreak{}4|\allowbreak{}0|\allowbreak{}-12|\allowbreak{}+++-\kern0pt{}-\kern0pt{}-+++-\kern0pt{}-\kern0pt{}- & \textit{Nombre de cellules}: 1\par \textit{Cellules (identifiants)}: 12\par \textit{Multisections}: quark\_\allowbreak{}down\_\allowbreak{}G3\par \textit{Collections}: quark\_\allowbreak{}down\par \textit{Couleur}: B\par \textit{Profils conjoints}: 1 & (80,\allowbreak{}66,\allowbreak{}4)\allowbreak{}x1 & (80,\allowbreak{}54,\allowbreak{}6)\allowbreak{}x1 \\
% VI-CATALOG-ROW familias 25
25\par 28|\allowbreak{}6|\allowbreak{}8|\allowbreak{}0|\allowbreak{}-12|\allowbreak{}+++-\kern0pt{}-\kern0pt{}-+++-\kern0pt{}-\kern0pt{}- & \textit{Nombre de cellules}: 1\par \textit{Cellules (identifiants)}: 62\par \textit{Multisections}: neutrino\_\allowbreak{}G3\par \textit{Collections}: neutrino\par \textit{Couleur}: none\par \textit{Profils conjoints}: 1 & (40,\allowbreak{}24,\allowbreak{}2)\allowbreak{}x1 & (80,\allowbreak{}54,\allowbreak{}6)\allowbreak{}x1 \\
% VI-CATALOG-ROW familias 26
26\par 30|\allowbreak{}-6|\allowbreak{}-4|\allowbreak{}-2|\allowbreak{}-8|\allowbreak{}0++-\kern0pt{}-\kern0pt{}-0++-\kern0pt{}-\kern0pt{}- & \textit{Nombre de cellules}: 1\par \textit{Cellules (identifiants)}: 4\par \textit{Multisections}: neutrino\_\allowbreak{}G4\par \textit{Collections}: neutrino\par \textit{Couleur}: none\par \textit{Profils conjoints}: 1 & (80,\allowbreak{}54,\allowbreak{}7)\allowbreak{}x1 & (56,\allowbreak{}54,\allowbreak{}5)\allowbreak{}x1 \\
% VI-CATALOG-ROW familias 27
27\par 30|\allowbreak{}-6|\allowbreak{}-8|\allowbreak{}0|\allowbreak{}-8|\allowbreak{}+0+-\kern0pt{}-\kern0pt{}-+0+-\kern0pt{}-\kern0pt{}- & \textit{Nombre de cellules}: 1\par \textit{Cellules (identifiants)}: 55\par \textit{Multisections}: charged\_\allowbreak{}lepton\_\allowbreak{}G4\par \textit{Collections}: charged\_\allowbreak{}lepton\par \textit{Couleur}: none\par \textit{Profils conjoints}: 1 & (60,\allowbreak{}78,\allowbreak{}25)\allowbreak{}x1 & (80,\allowbreak{}54,\allowbreak{}6)\allowbreak{}x1 \\
% VI-CATALOG-ROW familias 28
28\par 30|\allowbreak{}-6|\allowbreak{}0|\allowbreak{}-2|\allowbreak{}-8|\allowbreak{}++0-\kern0pt{}-\kern0pt{}-++0-\kern0pt{}-\kern0pt{}- & \textit{Nombre de cellules}: 2\par \textit{Cellules (identifiants)}: 24,\allowbreak{}34\par \textit{Multisections}: neutrino\_\allowbreak{}G2 | quark\_\allowbreak{}down\_\allowbreak{}G2\par \textit{Collections}: neutrino | quark\_\allowbreak{}down\par \textit{Couleur}: R | none\par \textit{Profils conjoints}: 2 & (0,\allowbreak{}24,\allowbreak{}0)\allowbreak{}x1; (60,\allowbreak{}78,\allowbreak{}25)\allowbreak{}x1 & (4,\allowbreak{}54,\allowbreak{}2)\allowbreak{}x1; (80,\allowbreak{}54,\allowbreak{}6)\allowbreak{}x1 \\
% VI-CATALOG-ROW familias 29
29\par 30|\allowbreak{}-6|\allowbreak{}0|\allowbreak{}-2|\allowbreak{}-8|\allowbreak{}0++-\kern0pt{}-\kern0pt{}-0++-\kern0pt{}-\kern0pt{}- & \textit{Nombre de cellules}: 1\par \textit{Cellules (identifiants)}: 75\par \textit{Multisections}: charged\_\allowbreak{}lepton\_\allowbreak{}G4\par \textit{Collections}: charged\_\allowbreak{}lepton\par \textit{Couleur}: none\par \textit{Profils conjoints}: 1 & (40,\allowbreak{}54,\allowbreak{}6)\allowbreak{}x1 & (56,\allowbreak{}54,\allowbreak{}5)\allowbreak{}x1 \\
% VI-CATALOG-ROW familias 30
30\par 30|\allowbreak{}-6|\allowbreak{}0|\allowbreak{}0|\allowbreak{}-12|\allowbreak{}+++-\kern0pt{}-\kern0pt{}-+++-\kern0pt{}-\kern0pt{}- & \textit{Nombre de cellules}: 1\par \textit{Cellules (identifiants)}: 14\par \textit{Multisections}: quark\_\allowbreak{}down\_\allowbreak{}G3\par \textit{Collections}: quark\_\allowbreak{}down\par \textit{Couleur}: R\par \textit{Profils conjoints}: 1 & (80,\allowbreak{}54,\allowbreak{}6)\allowbreak{}x1 & (80,\allowbreak{}54,\allowbreak{}6)\allowbreak{}x1 \\
% VI-CATALOG-ROW familias 31
31\par 30|\allowbreak{}-6|\allowbreak{}0|\allowbreak{}0|\allowbreak{}-8|\allowbreak{}+0+-\kern0pt{}-\kern0pt{}-+0+-\kern0pt{}-\kern0pt{}- & \textit{Nombre de cellules}: 1\par \textit{Cellules (identifiants)}: 54\par \textit{Multisections}: quark\_\allowbreak{}down\_\allowbreak{}G4\par \textit{Collections}: quark\_\allowbreak{}down\par \textit{Couleur}: R\par \textit{Profils conjoints}: 1 & (0,\allowbreak{}24,\allowbreak{}0)\allowbreak{}x1 & (80,\allowbreak{}54,\allowbreak{}6)\allowbreak{}x1 \\
% VI-CATALOG-ROW familias 32
32\par 30|\allowbreak{}-6|\allowbreak{}4|\allowbreak{}0|\allowbreak{}-12|\allowbreak{}+++-\kern0pt{}-\kern0pt{}-+++-\kern0pt{}-\kern0pt{}- & \textit{Nombre de cellules}: 1\par \textit{Cellules (identifiants)}: 65\par \textit{Multisections}: quark\_\allowbreak{}up\_\allowbreak{}G3\par \textit{Collections}: quark\_\allowbreak{}up\par \textit{Couleur}: R\par \textit{Profils conjoints}: 1 & (80,\allowbreak{}54,\allowbreak{}6)\allowbreak{}x1 & (80,\allowbreak{}54,\allowbreak{}6)\allowbreak{}x1 \\
% VI-CATALOG-ROW familias 33
33\par 30|\allowbreak{}-6|\allowbreak{}8|\allowbreak{}0|\allowbreak{}-12|\allowbreak{}+++-\kern0pt{}-\kern0pt{}-+++-\kern0pt{}-\kern0pt{}- & \textit{Nombre de cellules}: 1\par \textit{Cellules (identifiants)}: 44\par \textit{Multisections}: neutrino\_\allowbreak{}G3\par \textit{Collections}: neutrino\par \textit{Couleur}: none\par \textit{Profils conjoints}: 1 & (60,\allowbreak{}84,\allowbreak{}28)\allowbreak{}x1 & (80,\allowbreak{}54,\allowbreak{}6)\allowbreak{}x1 \\
% VI-CATALOG-ROW familias 34
34\par 30|\allowbreak{}6|\allowbreak{}-4|\allowbreak{}0|\allowbreak{}-12|\allowbreak{}+++-\kern0pt{}-\kern0pt{}-+++-\kern0pt{}-\kern0pt{}- & \textit{Nombre de cellules}: 1\par \textit{Cellules (identifiants)}: 68\par \textit{Multisections}: quark\_\allowbreak{}down\_\allowbreak{}G3\par \textit{Collections}: quark\_\allowbreak{}down\par \textit{Couleur}: R\par \textit{Profils conjoints}: 1 & (60,\allowbreak{}84,\allowbreak{}28)\allowbreak{}x1 & (80,\allowbreak{}54,\allowbreak{}6)\allowbreak{}x1 \\
% VI-CATALOG-ROW familias 35
35\par 30|\allowbreak{}6|\allowbreak{}-8|\allowbreak{}0|\allowbreak{}-12|\allowbreak{}+++-\kern0pt{}-\kern0pt{}-+++-\kern0pt{}-\kern0pt{}- & \textit{Nombre de cellules}: 1\par \textit{Cellules (identifiants)}: 17\par \textit{Multisections}: quark\_\allowbreak{}up\_\allowbreak{}G3\par \textit{Collections}: quark\_\allowbreak{}up\par \textit{Couleur}: R\par \textit{Profils conjoints}: 1 & (80,\allowbreak{}54,\allowbreak{}6)\allowbreak{}x1 & (80,\allowbreak{}54,\allowbreak{}6)\allowbreak{}x1 \\
% VI-CATALOG-ROW familias 36
36\par 30|\allowbreak{}6|\allowbreak{}0|\allowbreak{}0|\allowbreak{}-12|\allowbreak{}+++-\kern0pt{}-\kern0pt{}-+++-\kern0pt{}-\kern0pt{}- & \textit{Nombre de cellules}: 2\par \textit{Cellules (identifiants)}: 38,\allowbreak{}58\par \textit{Multisections}: neutrino\_\allowbreak{}G2 | neutrino\_\allowbreak{}G3\par \textit{Collections}: neutrino\par \textit{Couleur}: none\par \textit{Profils conjoints}: 2 & (60,\allowbreak{}78,\allowbreak{}25)\allowbreak{}x1; (80,\allowbreak{}54,\allowbreak{}6)\allowbreak{}x1 & (80,\allowbreak{}54,\allowbreak{}6)\allowbreak{}x2 \\
% VI-CATALOG-ROW familias 37
37\par 30|\allowbreak{}6|\allowbreak{}0|\allowbreak{}0|\allowbreak{}-8|\allowbreak{}+++-0-+++-0- & \textit{Nombre de cellules}: 1\par \textit{Cellules (identifiants)}: 78\par \textit{Multisections}: neutrino\_\allowbreak{}G4\par \textit{Collections}: neutrino\par \textit{Couleur}: none\par \textit{Profils conjoints}: 1 & (0,\allowbreak{}24,\allowbreak{}0)\allowbreak{}x1 & (80,\allowbreak{}54,\allowbreak{}6)\allowbreak{}x1 \\
% VI-CATALOG-ROW familias 38
38\par 30|\allowbreak{}6|\allowbreak{}0|\allowbreak{}2|\allowbreak{}-8|\allowbreak{}+++-\kern0pt{}-0+++-\kern0pt{}-0 & \textit{Nombre de cellules}: 1\par \textit{Cellules (identifiants)}: 27\par \textit{Multisections}: charged\_\allowbreak{}lepton\_\allowbreak{}G4\par \textit{Collections}: charged\_\allowbreak{}lepton\par \textit{Couleur}: none\par \textit{Profils conjoints}: 1 & (40,\allowbreak{}54,\allowbreak{}6)\allowbreak{}x1 & (56,\allowbreak{}54,\allowbreak{}5)\allowbreak{}x1 \\
% VI-CATALOG-ROW familias 39
39\par 30|\allowbreak{}6|\allowbreak{}0|\allowbreak{}2|\allowbreak{}-8|\allowbreak{}+++0-\kern0pt{}-+++0-\kern0pt{}- & \textit{Nombre de cellules}: 1\par \textit{Cellules (identifiants)}: 48\par \textit{Multisections}: quark\_\allowbreak{}down\_\allowbreak{}G2\par \textit{Collections}: quark\_\allowbreak{}down\par \textit{Couleur}: R\par \textit{Profils conjoints}: 1 & (0,\allowbreak{}24,\allowbreak{}0)\allowbreak{}x1 & (4,\allowbreak{}54,\allowbreak{}2)\allowbreak{}x1 \\
% VI-CATALOG-ROW familias 40
40\par 30|\allowbreak{}6|\allowbreak{}4|\allowbreak{}0|\allowbreak{}-12|\allowbreak{}+++-\kern0pt{}-\kern0pt{}-+++-\kern0pt{}-\kern0pt{}- & \textit{Nombre de cellules}: 1\par \textit{Cellules (identifiants)}: 7\par \textit{Multisections}: charged\_\allowbreak{}lepton\_\allowbreak{}G4\par \textit{Collections}: charged\_\allowbreak{}lepton\par \textit{Couleur}: none\par \textit{Profils conjoints}: 1 & (60,\allowbreak{}78,\allowbreak{}25)\allowbreak{}x1 & (80,\allowbreak{}54,\allowbreak{}6)\allowbreak{}x1 \\
% VI-CATALOG-ROW familias 41
41\par 30|\allowbreak{}6|\allowbreak{}8|\allowbreak{}0|\allowbreak{}-12|\allowbreak{}+++-\kern0pt{}-\kern0pt{}-+++-\kern0pt{}-\kern0pt{}- & \textit{Nombre de cellules}: 1\par \textit{Cellules (identifiants)}: 28\par \textit{Multisections}: quark\_\allowbreak{}down\_\allowbreak{}G4\par \textit{Collections}: quark\_\allowbreak{}down\par \textit{Couleur}: R\par \textit{Profils conjoints}: 1 & (80,\allowbreak{}54,\allowbreak{}7)\allowbreak{}x1 & (80,\allowbreak{}54,\allowbreak{}6)\allowbreak{}x1 \\
% VI-CATALOG-ROW familias 42
42\par 34|\allowbreak{}0|\allowbreak{}-12|\allowbreak{}-2|\allowbreak{}-8|\allowbreak{}0++-\kern0pt{}-\kern0pt{}-0++-\kern0pt{}-\kern0pt{}- & \textit{Nombre de cellules}: 1\par \textit{Cellules (identifiants)}: 45\par \textit{Multisections}: charged\_\allowbreak{}lepton\_\allowbreak{}G4\par \textit{Collections}: charged\_\allowbreak{}lepton\par \textit{Couleur}: none\par \textit{Profils conjoints}: 1 & (80,\allowbreak{}66,\allowbreak{}4)\allowbreak{}x1 & (56,\allowbreak{}54,\allowbreak{}5)\allowbreak{}x1 \\
% VI-CATALOG-ROW familias 43
43\par 34|\allowbreak{}0|\allowbreak{}-12|\allowbreak{}0|\allowbreak{}-12|\allowbreak{}+++-\kern0pt{}-\kern0pt{}-+++-\kern0pt{}-\kern0pt{}- & \textit{Nombre de cellules}: 2\par \textit{Cellules (identifiants)}: 57,\allowbreak{}76\par \textit{Multisections}: charged\_\allowbreak{}lepton\_\allowbreak{}G2 | neutrino\_\allowbreak{}G4\par \textit{Collections}: charged\_\allowbreak{}lepton | neutrino\par \textit{Couleur}: none\par \textit{Profils conjoints}: 2 & (40,\allowbreak{}84,\allowbreak{}30)\allowbreak{}x1; (80,\allowbreak{}66,\allowbreak{}2)\allowbreak{}x1 & (80,\allowbreak{}54,\allowbreak{}6)\allowbreak{}x2 \\
% VI-CATALOG-ROW familias 44
44\par 34|\allowbreak{}0|\allowbreak{}-16|\allowbreak{}0|\allowbreak{}-8|\allowbreak{}+++-0-+++-0- & \textit{Nombre de cellules}: 1\par \textit{Cellules (identifiants)}: 37\par \textit{Multisections}: charged\_\allowbreak{}lepton\_\allowbreak{}G4\par \textit{Collections}: charged\_\allowbreak{}lepton\par \textit{Couleur}: none\par \textit{Profils conjoints}: 1 & (40,\allowbreak{}24,\allowbreak{}2)\allowbreak{}x1 & (80,\allowbreak{}54,\allowbreak{}6)\allowbreak{}x1 \\
% VI-CATALOG-ROW familias 45
45\par 34|\allowbreak{}0|\allowbreak{}-16|\allowbreak{}0|\allowbreak{}-8|\allowbreak{}+0+-\kern0pt{}-\kern0pt{}-+0+-\kern0pt{}-\kern0pt{}- & \textit{Nombre de cellules}: 1\par \textit{Cellules (identifiants)}: 15\par \textit{Multisections}: quark\_\allowbreak{}up\_\allowbreak{}G3\par \textit{Collections}: quark\_\allowbreak{}up\par \textit{Couleur}: B\par \textit{Profils conjoints}: 1 & (40,\allowbreak{}78,\allowbreak{}27)\allowbreak{}x1 & (80,\allowbreak{}54,\allowbreak{}6)\allowbreak{}x1 \\
% VI-CATALOG-ROW familias 46
46\par 34|\allowbreak{}0|\allowbreak{}-20|\allowbreak{}-2|\allowbreak{}-8|\allowbreak{}0++-\kern0pt{}-\kern0pt{}-0++-\kern0pt{}-\kern0pt{}- & \textit{Nombre de cellules}: 1\par \textit{Cellules (identifiants)}: 35\par \textit{Multisections}: quark\_\allowbreak{}up\_\allowbreak{}G3\par \textit{Collections}: quark\_\allowbreak{}up\par \textit{Couleur}: B\par \textit{Profils conjoints}: 1 & (100,\allowbreak{}66,\allowbreak{}0)\allowbreak{}x1 & (56,\allowbreak{}54,\allowbreak{}5)\allowbreak{}x1 \\
% VI-CATALOG-ROW familias 47
47\par 34|\allowbreak{}0|\allowbreak{}-4|\allowbreak{}-2|\allowbreak{}-8|\allowbreak{}++0-\kern0pt{}-\kern0pt{}-++0-\kern0pt{}-\kern0pt{}- & \textit{Nombre de cellules}: 1\par \textit{Cellules (identifiants)}: 56\par \textit{Multisections}: neutrino\_\allowbreak{}G3\par \textit{Collections}: neutrino\par \textit{Couleur}: none\par \textit{Profils conjoints}: 1 & (40,\allowbreak{}24,\allowbreak{}2)\allowbreak{}x1 & (36,\allowbreak{}54,\allowbreak{}4)\allowbreak{}x1 \\
% VI-CATALOG-ROW familias 48
48\par 34|\allowbreak{}0|\allowbreak{}-4|\allowbreak{}0|\allowbreak{}-12|\allowbreak{}+++-\kern0pt{}-\kern0pt{}-+++-\kern0pt{}-\kern0pt{}- & \textit{Nombre de cellules}: 3\par \textit{Cellules (identifiants)}: 6,\allowbreak{}25,\allowbreak{}46\par \textit{Multisections}: charged\_\allowbreak{}lepton\_\allowbreak{}G2 | neutrino\_\allowbreak{}G4 | quark\_\allowbreak{}down\_\allowbreak{}G4\par \textit{Collections}: charged\_\allowbreak{}lepton | neutrino | quark\_\allowbreak{}down\par \textit{Couleur}: B | none\par \textit{Profils conjoints}: 2 & (80,\allowbreak{}66,\allowbreak{}2)\allowbreak{}x1; (80,\allowbreak{}66,\allowbreak{}3)\allowbreak{}x2 & (80,\allowbreak{}54,\allowbreak{}6)\allowbreak{}x3 \\
% VI-CATALOG-ROW familias 49
49\par 34|\allowbreak{}0|\allowbreak{}-4|\allowbreak{}2|\allowbreak{}-8|\allowbreak{}+++-\kern0pt{}-0+++-\kern0pt{}-0 & \textit{Nombre de cellules}: 1\par \textit{Cellules (identifiants)}: 67\par \textit{Multisections}: quark\_\allowbreak{}up\_\allowbreak{}G3\par \textit{Collections}: quark\_\allowbreak{}up\par \textit{Couleur}: G\par \textit{Profils conjoints}: 1 & (100,\allowbreak{}66,\allowbreak{}2)\allowbreak{}x1 & (56,\allowbreak{}54,\allowbreak{}5)\allowbreak{}x1 \\
% VI-CATALOG-ROW familias 50
50\par 34|\allowbreak{}0|\allowbreak{}-8|\allowbreak{}-2|\allowbreak{}-8|\allowbreak{}++0-\kern0pt{}-\kern0pt{}-++0-\kern0pt{}-\kern0pt{}- & \textit{Nombre de cellules}: 1\par \textit{Cellules (identifiants)}: 66\par \textit{Multisections}: quark\_\allowbreak{}down\_\allowbreak{}G3\par \textit{Collections}: quark\_\allowbreak{}down\par \textit{Couleur}: B\par \textit{Profils conjoints}: 1 & (40,\allowbreak{}78,\allowbreak{}27)\allowbreak{}x1 & (80,\allowbreak{}54,\allowbreak{}6)\allowbreak{}x1 \\
% VI-CATALOG-ROW familias 51
51\par 34|\allowbreak{}0|\allowbreak{}-8|\allowbreak{}-2|\allowbreak{}-8|\allowbreak{}0++-\kern0pt{}-\kern0pt{}-0++-\kern0pt{}-\kern0pt{}- & \textit{Nombre de cellules}: 1\par \textit{Cellules (identifiants)}: 77\par \textit{Multisections}: charged\_\allowbreak{}lepton\_\allowbreak{}G4\par \textit{Collections}: charged\_\allowbreak{}lepton\par \textit{Couleur}: none\par \textit{Profils conjoints}: 1 & (80,\allowbreak{}66,\allowbreak{}3)\allowbreak{}x1 & (56,\allowbreak{}54,\allowbreak{}5)\allowbreak{}x1 \\
% VI-CATALOG-ROW familias 52
52\par 34|\allowbreak{}0|\allowbreak{}-8|\allowbreak{}0|\allowbreak{}-12|\allowbreak{}+++-\kern0pt{}-\kern0pt{}-+++-\kern0pt{}-\kern0pt{}- & \textit{Nombre de cellules}: 1\par \textit{Cellules (identifiants)}: 36\par \textit{Multisections}: quark\_\allowbreak{}down\_\allowbreak{}G4\par \textit{Collections}: quark\_\allowbreak{}down\par \textit{Couleur}: G\par \textit{Profils conjoints}: 1 & (40,\allowbreak{}84,\allowbreak{}30)\allowbreak{}x1 & (80,\allowbreak{}54,\allowbreak{}6)\allowbreak{}x1 \\
% VI-CATALOG-ROW familias 53
53\par 34|\allowbreak{}0|\allowbreak{}-8|\allowbreak{}0|\allowbreak{}-8|\allowbreak{}+0+-\kern0pt{}-\kern0pt{}-+0+-\kern0pt{}-\kern0pt{}- & \textit{Nombre de cellules}: 1\par \textit{Cellules (identifiants)}: 5\par \textit{Multisections}: charged\_\allowbreak{}lepton\_\allowbreak{}G4\par \textit{Collections}: charged\_\allowbreak{}lepton\par \textit{Couleur}: none\par \textit{Profils conjoints}: 1 & (40,\allowbreak{}24,\allowbreak{}2)\allowbreak{}x1 & (80,\allowbreak{}54,\allowbreak{}6)\allowbreak{}x1 \\
% VI-CATALOG-ROW familias 54
54\par 34|\allowbreak{}0|\allowbreak{}0|\allowbreak{}-2|\allowbreak{}-8|\allowbreak{}++0-\kern0pt{}-\kern0pt{}-++0-\kern0pt{}-\kern0pt{}- & \textit{Nombre de cellules}: 1\par \textit{Cellules (identifiants)}: 26\par \textit{Multisections}: neutrino\_\allowbreak{}G3\par \textit{Collections}: neutrino\par \textit{Couleur}: none\par \textit{Profils conjoints}: 1 & (40,\allowbreak{}78,\allowbreak{}27)\allowbreak{}x1 & (80,\allowbreak{}54,\allowbreak{}6)\allowbreak{}x1 \\
% VI-CATALOG-ROW familias 55
55\par 34|\allowbreak{}0|\allowbreak{}0|\allowbreak{}0|\allowbreak{}-8|\allowbreak{}+0+-\kern0pt{}-\kern0pt{}-+0+-\kern0pt{}-\kern0pt{}- & \textit{Nombre de cellules}: 1\par \textit{Cellules (identifiants)}: 47\par \textit{Multisections}: quark\_\allowbreak{}up\_\allowbreak{}G3\par \textit{Collections}: quark\_\allowbreak{}up\par \textit{Couleur}: G\par \textit{Profils conjoints}: 1 & (40,\allowbreak{}78,\allowbreak{}27)\allowbreak{}x1 & (80,\allowbreak{}54,\allowbreak{}6)\allowbreak{}x1 \\
% VI-CATALOG-ROW familias 56
56\par 34|\allowbreak{}0|\allowbreak{}0|\allowbreak{}2|\allowbreak{}-8|\allowbreak{}+++0-\kern0pt{}-+++0-\kern0pt{}- & \textit{Nombre de cellules}: 1\par \textit{Cellules (identifiants)}: 16\par \textit{Multisections}: quark\_\allowbreak{}down\_\allowbreak{}G3\par \textit{Collections}: quark\_\allowbreak{}down\par \textit{Couleur}: G\par \textit{Profils conjoints}: 1 & (40,\allowbreak{}24,\allowbreak{}2)\allowbreak{}x1 & (36,\allowbreak{}54,\allowbreak{}4)\allowbreak{}x1 \\
\end{longtable}
\endgroup
\subsection{Classification des 23 secteurs documentés}

Ces lignes réunissent des secteurs explicites, des résidences et une représentation. La clé
d’état renvoie à la légende contiguë et conserve le statut de la source.
En particulier, un opérateur interne exact, une évaluation conditionnée et une
réalisation scalaire fermée sont maintenus comme des affirmations distinctes. Le tableau
ne met pas leur statut à jour et ne transforme pas un registre conditionné en prédiction.

\begingroup
\setlength{\tabcolsep}{4pt}
\setlength{\extrarowheight}{4pt}
\renewcommand{\arraystretch}{1.13}
\begin{longtable}{@{}>{\raggedright\arraybackslash}p{\dimexpr 0.08860759\linewidth-2\tabcolsep\relax}>{\raggedright\arraybackslash}p{\dimexpr 0.41772152\linewidth-2\tabcolsep\relax}>{\raggedright\arraybackslash}p{\dimexpr 0.49367089\linewidth-2\tabcolsep\relax}@{}}
\caption{États originaux des classes.}\label{tab:vi-estados-clases}\\
\toprule
Clé & État littéral & Lecture documentaire \\
\midrule
\endfirsthead
\multicolumn{3}{l}{\textit{Suite}}\\
\toprule
Clé & État littéral & Lecture documentaire \\
\midrule
\endhead
\midrule
\multicolumn{3}{r}{Suite à la page suivante}\\
\endfoot
\bottomrule
\endlastfoot
C01 & CLOSED\_\allowbreak{}ELECTRON\_\allowbreak{}TWO\_\allowbreak{}WAY & Électron : fermeture bidirectionnelle documentée. \\
C02 & CLOSED\_\allowbreak{}NULL\_\allowbreak{}CHART & Coordinatisation nulle fermée à cet état arrêté. \\
C03 & NOT\_\allowbreak{}A\_\allowbreak{}UNIVERSAL\_\allowbreak{}MASS\_\allowbreak{}SCALAR & La classe ne s’exprime pas comme un scalaire universel de masse. \\
C04 & OPEN\_\allowbreak{}FLAVOR\_\allowbreak{}COLOR\_\allowbreak{}SCALARIZATION & Scalarisation de saveur et de couleur ouverte à cet état arrêté. \\
C05 & OPEN\_\allowbreak{}MIXING\_\allowbreak{}AND\_\allowbreak{}SCALE & Mélange et échelle ouverts à cet état arrêté. \\
C06 & OPEN\_\allowbreak{}NAMED\_\allowbreak{}STATE\_\allowbreak{}SCALARIZATION & Scalarisation de l’état individualisé ouverte à cet état arrêté. \\
C07 & OPEN\_\allowbreak{}PHYSICAL\_\allowbreak{}SCALARIZATION & Scalarisation physique ouverte à cet état arrêté. \\
C08 & OPEN\_\allowbreak{}ROUTE\_\allowbreak{}AND\_\allowbreak{}SCALARIZATION & Route et scalarisation ouvertes à cet état arrêté. \\
\end{longtable}
\endgroup

\begingroup
\setlength{\tabcolsep}{4pt}
\setlength{\extrarowheight}{4pt}
\renewcommand{\arraystretch}{1.08}
\begin{longtable}{@{}>{\raggedright\arraybackslash}p{\dimexpr 0.17307692\linewidth-2\tabcolsep\relax}>{\raggedright\arraybackslash}p{\dimexpr 0.38461538\linewidth-2\tabcolsep\relax}>{\raggedright\arraybackslash}p{\dimexpr 0.33333333\linewidth-2\tabcolsep\relax}>{\raggedright\arraybackslash}p{\dimexpr 0.10897436\linewidth-2\tabcolsep\relax}@{}}
\caption{Les 23 classes du catalogue, dans leur ordre documentaire.}\label{tab:vi-clases}\\
\toprule
Classe & Représentants et résidence & Incidence et représentation & État \\
\midrule
\endfirsthead
\multicolumn{4}{l}{\textit{Suite}}\\
\toprule
Classe & Représentants et résidence & Incidence et représentation & État \\
\midrule
\endhead
\midrule
\multicolumn{4}{r}{Suite à la page suivante}\\
\endfoot
\bottomrule
\endlastfoot
% VI-CATALOG-ROW clases 1
SM-LC-1\par lepton\_\allowbreak{}cargado & \textit{Représentants}: e− / e+\par \textit{Résidence}: centre (5,\allowbreak{}5) et antisection & \textit{Cellules sélectionnées}: 1\par \textit{Multisections}: charged\_\allowbreak{}lepton\_\allowbreak{}G0\par \textit{Charge}: ±1 e\par \textit{Représentation}: sans couleur & C01 \\
% VI-CATALOG-ROW clases 2
SM-LC-2\par lepton\_\allowbreak{}cargado & \textit{Représentants}: μ− / μ+\par \textit{Résidence}: multisection de lepton chargé & \textit{Cellules sélectionnées}: 8\par \textit{Multisections}: charged\_\allowbreak{}lepton\_\allowbreak{}G2\par \textit{Charge}: ±1 e\par \textit{Représentation}: sans couleur & C07 \\
% VI-CATALOG-ROW clases 3
SM-LC-3\par lepton\_\allowbreak{}cargado & \textit{Représentants}: τ− / τ+\par \textit{Résidence}: multisection de lepton chargé & \textit{Cellules sélectionnées}: 16\par \textit{Multisections}: charged\_\allowbreak{}lepton\_\allowbreak{}G4\par \textit{Charge}: ±1 e\par \textit{Représentation}: sans couleur & C07 \\
% VI-CATALOG-ROW clases 4
SM-NU-1\par neutrino & \textit{Représentants}: ν\_\allowbreak{}e / anti-ν\_\allowbreak{}e\par \textit{Résidence}: multisections de neutrino & \textit{Cellules sélectionnées}: 20\par \textit{Multisections}: neutrino\_\allowbreak{}G1 | neutrino\_\allowbreak{}G2 | neutrino\_\allowbreak{}G3 | neutrino\_\allowbreak{}G4\par \textit{Charge}: 0\par \textit{Représentation}: sans couleur & C05 \\
% VI-CATALOG-ROW clases 5
SM-NU-2\par neutrino & \textit{Représentants}: ν\_\allowbreak{}μ / anti-ν\_\allowbreak{}μ\par \textit{Résidence}: multisections de neutrino & \textit{Cellules sélectionnées}: 20\par \textit{Multisections}: neutrino\_\allowbreak{}G1 | neutrino\_\allowbreak{}G2 | neutrino\_\allowbreak{}G3 | neutrino\_\allowbreak{}G4\par \textit{Charge}: 0\par \textit{Représentation}: sans couleur & C05 \\
% VI-CATALOG-ROW clases 6
SM-NU-3\par neutrino & \textit{Représentants}: ν\_\allowbreak{}τ / anti-ν\_\allowbreak{}τ\par \textit{Résidence}: multisections de neutrino & \textit{Cellules sélectionnées}: 20\par \textit{Multisections}: neutrino\_\allowbreak{}G1 | neutrino\_\allowbreak{}G2 | neutrino\_\allowbreak{}G3 | neutrino\_\allowbreak{}G4\par \textit{Charge}: 0\par \textit{Représentation}: sans couleur & C05 \\
% VI-CATALOG-ROW clases 7
SM-Q-U1\par quark & \textit{Représentants}: u / anti-u\par \textit{Résidence}: multisection de type haut & \textit{Cellules sélectionnées}: 16\par \textit{Multisections}: quark\_\allowbreak{}up\_\allowbreak{}G1 | quark\_\allowbreak{}up\_\allowbreak{}G3\par \textit{Charge}: ±2/\allowbreak{}3 e\par \textit{Représentation}: R,\allowbreak{}G,\allowbreak{}B = orbite Z/\allowbreak{}3 & C04 \\
% VI-CATALOG-ROW clases 8
SM-Q-D1\par quark & \textit{Représentants}: d / anti-d\par \textit{Résidence}: multisection de type bas & \textit{Cellules sélectionnées}: 20\par \textit{Multisections}: quark\_\allowbreak{}down\_\allowbreak{}G1 | quark\_\allowbreak{}down\_\allowbreak{}G2 | quark\_\allowbreak{}down\_\allowbreak{}G3 | quark\_\allowbreak{}down\_\allowbreak{}G4\par \textit{Charge}: ±1/\allowbreak{}3 e\par \textit{Représentation}: R,\allowbreak{}G,\allowbreak{}B = orbite Z/\allowbreak{}3 & C04 \\
% VI-CATALOG-ROW clases 9
SM-Q-U2\par quark & \textit{Représentants}: c / anti-c\par \textit{Résidence}: multisection de type haut & \textit{Cellules sélectionnées}: 16\par \textit{Multisections}: quark\_\allowbreak{}up\_\allowbreak{}G1 | quark\_\allowbreak{}up\_\allowbreak{}G3\par \textit{Charge}: ±2/\allowbreak{}3 e\par \textit{Représentation}: R,\allowbreak{}G,\allowbreak{}B = orbite Z/\allowbreak{}3 & C04 \\
% VI-CATALOG-ROW clases 10
SM-Q-D2\par quark & \textit{Représentants}: s / anti-s\par \textit{Résidence}: multisection de type bas & \textit{Cellules sélectionnées}: 20\par \textit{Multisections}: quark\_\allowbreak{}down\_\allowbreak{}G1 | quark\_\allowbreak{}down\_\allowbreak{}G2 | quark\_\allowbreak{}down\_\allowbreak{}G3 | quark\_\allowbreak{}down\_\allowbreak{}G4\par \textit{Charge}: ±1/\allowbreak{}3 e\par \textit{Représentation}: R,\allowbreak{}G,\allowbreak{}B = orbite Z/\allowbreak{}3 & C04 \\
% VI-CATALOG-ROW clases 11
SM-Q-U3\par quark & \textit{Représentants}: t / anti-t\par \textit{Résidence}: multisection de type haut & \textit{Cellules sélectionnées}: 16\par \textit{Multisections}: quark\_\allowbreak{}up\_\allowbreak{}G1 | quark\_\allowbreak{}up\_\allowbreak{}G3\par \textit{Charge}: ±2/\allowbreak{}3 e\par \textit{Représentation}: R,\allowbreak{}G,\allowbreak{}B = orbite Z/\allowbreak{}3 & C04 \\
% VI-CATALOG-ROW clases 12
SM-Q-D3\par quark & \textit{Représentants}: b / anti-b\par \textit{Résidence}: multisection de type bas & \textit{Cellules sélectionnées}: 20\par \textit{Multisections}: quark\_\allowbreak{}down\_\allowbreak{}G1 | quark\_\allowbreak{}down\_\allowbreak{}G2 | quark\_\allowbreak{}down\_\allowbreak{}G3 | quark\_\allowbreak{}down\_\allowbreak{}G4\par \textit{Charge}: ±1/\allowbreak{}3 e\par \textit{Représentation}: R,\allowbreak{}G,\allowbreak{}B = orbite Z/\allowbreak{}3 & C04 \\
% VI-CATALOG-ROW clases 13
SM-GA-EM\par boson\_\allowbreak{}gauge & \textit{Représentants}: γ\par \textit{Résidence}: secteur de jauge nul & \textit{Cellules sélectionnées}: 0\par \textit{Multisections}: \textemdash{}\par \textit{Charge}: 0\par \textit{Représentation}: sans couleur & C02 \\
% VI-CATALOG-ROW clases 14
SM-GA-GL\par boson\_\allowbreak{}gauge & \textit{Représentants}: g\par \textit{Résidence}: secteur de jauge de couleur & \textit{Cellules sélectionnées}: 0\par \textit{Multisections}: \textemdash{}\par \textit{Charge}: 0\par \textit{Représentation}: octet de jauge & C02 \\
% VI-CATALOG-ROW clases 15
SM-GA-W\par boson\_\allowbreak{}gauge & \textit{Représentants}: W+ / W−\par \textit{Résidence}: secteur de jauge massif & \textit{Cellules sélectionnées}: 0\par \textit{Multisections}: \textemdash{}\par \textit{Charge}: ±1 e\par \textit{Représentation}: sans couleur & C08 \\
% VI-CATALOG-ROW clases 16
SM-GA-Z\par boson\_\allowbreak{}gauge & \textit{Représentants}: Z0\par \textit{Résidence}: secteur de jauge massif & \textit{Cellules sélectionnées}: 0\par \textit{Multisections}: \textemdash{}\par \textit{Charge}: 0\par \textit{Représentation}: sans couleur & C08 \\
% VI-CATALOG-ROW clases 17
SM-H-1\par escalar & \textit{Représentants}: H\par \textit{Résidence}: secteur scalaire & \textit{Cellules sélectionnées}: 0\par \textit{Multisections}: \textemdash{}\par \textit{Charge}: 0\par \textit{Représentation}: sans couleur & C08 \\
% VI-CATALOG-ROW clases 18
COMP-MESON\par meson & \textit{Représentants}: π0, π±, K0, K±, η, ρ, φ, J/\allowbreak{}ψ, Υ…\par \textit{Résidence}: produit de section et d’antisection & \textit{Cellules sélectionnées}: 0\par \textit{Multisections}: \textemdash{}\par \textit{Charge}: dépendant de l’état composé\par \textit{Représentation}: singulet ou autres représentations & C06 \\
% VI-CATALOG-ROW clases 19
COMP-BARYON\par barion & \textit{Représentants}: p, n et configurations baryoniques\par \textit{Résidence}: produit de trois sections & \textit{Cellules sélectionnées}: 0\par \textit{Multisections}: \textemdash{}\par \textit{Charge}: dépendant de l’état composé\par \textit{Représentation}: singulet physique & C06 \\
% VI-CATALOG-ROW clases 20
TOPO-FIB\par anyon\_\allowbreak{}no\_\allowbreak{}abeliano & \textit{Représentants}: 1, τ\par \textit{Résidence}: secteur topologique & \textit{Cellules sélectionnées}: 0\par \textit{Multisections}: \textemdash{}\par \textit{Charge}: définie par le codomaine topologique\par \textit{Représentation}: sans objet & C03 \\
% VI-CATALOG-ROW clases 21
TOPO-ISING\par sector\_\allowbreak{}majorana & \textit{Représentants}: 1, ψ, σ\par \textit{Résidence}: secteur topologique / zéro & \textit{Cellules sélectionnées}: 0\par \textit{Multisections}: \textemdash{}\par \textit{Charge}: dépendant de la réalisation\par \textit{Représentation}: sans couleur & C03 \\
% VI-CATALOG-ROW clases 22
TOPO-Z6\par anyones\_\allowbreak{}abelianos & \textit{Représentants}: boson, fermion et quatre secteurs tressés\par \textit{Résidence}: secteur topologique Z/\allowbreak{}6 & \textit{Cellules sélectionnées}: 0\par \textit{Multisections}: \textemdash{}\par \textit{Charge}: Z/\allowbreak{}6 définie par le tressage\par \textit{Représentation}: sans couleur & C03 \\
% VI-CATALOG-ROW clases 23
VIRT-1\par virtual & \textit{Représentants}: sections virtuelles\par \textit{Résidence}: système projectif TPK & \textit{Cellules sélectionnées}: 0\par \textit{Multisections}: \textemdash{}\par \textit{Charge}: dépendant du secteur de prolongement\par \textit{Représentation}: selon le secteur & C03 \\
\end{longtable}
\endgroup
